\documentclass[12pt]{scrartcl}
\usepackage[utf8]{inputenc}
\usepackage[T1]{fontenc}
\usepackage[english,ngerman]{babel}
\usepackage[sorting=none,style=ieee]{biblatex}
\usepackage[margin=25mm]{geometry}
\usepackage{scrlayer-scrpage,anyfontsize,hyperref}

\renewcommand{\familydefault}{\sfdefault}
\setkomafont{section}{\fontsize{14}{14}\selectfont}
\usepackage{fontspec}
\setmainfont[
 Ligatures=TeX,
 BoldFont={CALIBRIB.TTF}, 
 ItalicFont={CALIBRII.TTF}
 ]{CALIBRI.TTF}

\ihead{Titel}
\ohead{Autor}
\cfoot{}
\ofoot{\thepage}

\bibliography{quellen}

\begin{document}

\begin{center}
\textbf{\fontsize{16}{16}\selectfont Kurzfassung Proseminararbeit}
\end{center}

\bigskip

\section*{Motivation}

\textit{ausstehend} Warum interessant? Warum lohnt sich Zuhören? Beispiel zur Verdeutlichung der Problemstellung? Was lernen wir?

\section*{Grundlagen und Voraussetzungen}

\textit{ausstehend} Was muss ich wissen? Was sind technische Voraussetzungen? Wie sieht konkretes Problem aus? Case Studies: In welchem Umfeld?

\section*{Vorgehensweisen und Modelle}

\textit{ausstehend} Weg zum Ziel? Problemlösung? Modelle? Experimente?

\section*{Erkenntnisse}

\textit{ausstehend} Was haben wir gelernt? Verallgemeinerung? Konkrete Problemlösung? Warum ist sie richtig?

\section*{Kritik und Ausblick}

\textit{ausstehend} Mehr Dimensionen? Vollständigkeit? Offene Punkte? Weitergehende Forschung?

Zitate benutzen! \cite{microBpm}
\section*{Quellen}
\printbibliography[heading=none]

\end{document}